\section{Discussion}

Our results suggest that explicit orchestration deserves to be studied as a first-class systems problem for LLM agents. The current literature often emphasizes stronger tools, larger models, or more sophisticated prompting strategies. In contrast, our study highlights that the \emph{control layer} between the model and the environment also plays an important role. Even lightweight utility design can significantly influence how often the agent acts, when it stops, and how much intermediate computation it accumulates.

At the same time, this work has several important limitations. First, the utility components are heuristic rather than learned. Both estimated gain and uncertainty are self-estimated signals, not calibrated probabilities. Second, the current policy is not shown to universally outperform stronger free-form baselines such as ReAct. Our contribution is therefore better understood as a structured orchestration framework and empirical analysis, rather than a universally best-performing method. Third, the current redundancy mechanism improves token compactness more clearly than latency, indicating that better semantic control mechanisms remain an open direction.

These limitations also point to future work. One natural extension is to replace some heuristic signals with learned scorers or calibrated uncertainty estimates. Another is to extend the framework to more complex environments and longer-horizon tasks. It would also be valuable to investigate explicit memory-aware orchestration, where retrieval, memory reading, and memory writing are controlled under the same decision framework. More broadly, our findings support the view that efficient LLM agents require not only strong reasoning models, but also explicit decision policies that constrain and organize multi-step behavior.